\documentclass[11pt,a4paper]{article}

% ------------------------------------------------------------------
% Setup (compile with pdflatex from this folder; figures/ and results/tables/ are picked up automatically)
% ------------------------------------------------------------------
\usepackage[margin=2.5cm]{geometry}
\usepackage[T1]{fontenc}
\usepackage{lmodern}
\usepackage{amsmath,amssymb,bm}
\usepackage{graphicx}
\usepackage{booktabs,array}
\usepackage{float}
\usepackage{placeins}
\usepackage{caption}
\usepackage{enumitem}
\usepackage{setspace}
\usepackage[hidelinks]{hyperref}
\onehalfspacing
\emergencystretch=3em
\captionsetup{font=small,labelfont=bf}
\setlist{itemsep=2pt,topsep=4pt}

\newcommand{\Tout}{T_{\mathrm{out,ele}}}
\newcommand{\Tin}{T_{\mathrm{in,ele}}}
\newcommand{\Tc}{T_{\mathrm{out,c}}}
\newcommand{\Vcell}{V_{\mathrm{cell}}}
\newcommand{\dT}{\Delta T}

% \fig{file-stem}{width-fraction}{caption}{label}: figure from figures/<file-stem>.png
\newcommand{\fig}[4]{%
  \begin{figure}[H]\centering
  \includegraphics[width=#2\textwidth]{figures/#1.png}
  \caption{#3}\label{#4}
  \end{figure}}

% \tab{name}: the table body written by the code, results/tables/<name>.tex
\newcommand{\tab}[1]{\input{results/tables/#1.tex}}

\title{\vspace*{2cm}\textbf{AI-Based Steady-State and Dynamic Modeling of the Thermal Behavior of a PEM Water Electrolyzer}\\[6pt]
\large Course Project --- Stage 1 Report\\[2pt]
\large CHL 4380: Applications of AI in Chemical Engineering}
\author{Jatin\\ Roll No.\ B23CH1022\\ Department of Chemical Engineering\\ Indian Institute of Technology Jodhpur\\[8pt]
Instructor: Dr.\ Varanasi Santhosh Kumar}
\date{}

\begin{document}
\maketitle
\thispagestyle{empty}
\vspace{1.5cm}

% ==================================================================
\begin{abstract}
\noindent
This report covers Stage~1 (modeling) of the course project. The process is the thermal management loop of a PEM water electrolyzer: a nonlinear, multi-input, disturbance-driven system. A three-state first-principles model, based on Zerrougui et al.\ (2026), is used as a data generator. Steady-state data come from a Latin-hypercube design of the stack current and two valve openings; dynamic data come from random multi-level excitation integrated with RK4. The steady-state map is learned with multiple linear regression, polynomial ridge regression, an MLP and a Takagi--Sugeno (ANFIS-type) model. The dynamics are learned with ARX, a NARX-MLP, an LSTM and a self-loop physics-informed network (PINC) trained without next-state targets. All models are compared on held-out data. The best steady-state model (MLP) reaches a test RMSE of 0.036~K on $\Tin$ and 0.00044~V on $\Vcell$. The best dynamic model (NARX-MLP) reaches a 2000~s free-run rollout RMSE of 0.012~K. All data are simulated, and several simulator parameters are assumed.
\end{abstract}
\clearpage

\tableofcontents
\clearpage

% ==================================================================
\section{Introduction}
PEM electrolyzers produce hydrogen from renewable electricity and respond quickly to load changes \cite{zerrougui2026,carmo2013}. Because the cell voltage exceeds the thermoneutral voltage, the stack is exothermic and needs a cooling loop. Poor thermal management lowers efficiency and accelerates degradation. A first-principles thermal model is awkward inside an optimizer or a model predictive controller, because it has to be integrated numerically and depends on uncertain parameters. A fast AI surrogate trained on simulated data avoids this \cite{lectures}.

\textbf{Objectives of Stage~1.}
\begin{enumerate}[label=\textbf{O\arabic*.}]
  \item Identify and justify the inputs and outputs (rubric a, d).
  \item Generate steady-state and dynamic datasets from a verified simulator (b).
  \item Preprocess the data without leakage (c).
  \item Train, validate and test steady-state and dynamic AI models, with simple baselines (e).
  \item Evaluate accuracy, rollout stability, cost and noise robustness (f).
\end{enumerate}
All data are generated by simulation, so every accuracy figure is a surrogate-versus-simulator comparison. Figure~\ref{fig:workflow} shows the workflow.

\fig{fig_workflow}{0.95}{Stage~1 workflow: first-principles model, simulated experiments, preprocessing and splitting, training and validation, test evaluation, and selected surrogates for Stages~2 and~3.}{fig:workflow}

% ==================================================================
\section{Process and First-Principles Model}
\subsection{Process description}
Cells are stacked in series ($n_{\mathrm{cell}}=26$). When the cell voltage $\Vcell$ exceeds $V_{\mathrm{th}}=1.48$~V, the excess is released as heat, so heat generation grows with the stack current $I$, which is an externally imposed disturbance. As shown in Fig.~\ref{fig:schematic}, a hot water loop (stack, separator, pump) exchanges heat with a cold coolant loop (pump, chiller) through a heat exchanger. The manipulated variables are the two valve openings $u_1$ (hot side) and $u_2$ (cold side) in $[0,1]$. The long-term goal of the project is to control $\Tin$ and $\dT=\Tout-\Tin$.

\fig{fig_system_schematic}{0.8}{Hot water loop and cold coolant loop of the PEM electrolyzer, with the measured temperatures, the valve openings $u_1,u_2$ and the current $I$.}{fig:schematic}

The process suits AI modeling because it is nonlinear (the heat generation depends on temperature and current, and cooling terms are products $u_i x_j$), multivariable (two valves and one disturbance act on three coupled temperatures), and is driven by a strongly fluctuating disturbance.

\subsection{State equations}
Assumptions: lumped capacitance for each of three nodes, separator temperature equal to $\Tout$, flows proportional to valve openings ($\dot m_{\mathrm{hot}}=u_1\dot m_0$, $\dot m_{\mathrm{cold}}=u_2\dot m_0$), constant fluid properties, ambient and chiller temperatures, and no gas holdup or degradation. With states $x_1=\Tout$, $x_2=\Tin$, $x_3=\Tc$,
\begin{align}
\dot x_1 &= \frac{1}{C_{\mathrm{th}}}\Big[\, n_{\mathrm{cell}} I\big(\Vcell(x_1,I)-V_{\mathrm{th}}\big) - hA_s\,(x_1-T_{\mathrm{amb}}) - u_1\dot m_0 c_p\,(x_1-x_2)\Big],\\
\dot x_2 &= \frac{1}{C_{h}}\Big[\, u_1\dot m_0 c_p\,(x_1-x_2) + UA_{\mathrm{ex}}\,(x_3-x_2)\Big],\\
\dot x_3 &= \frac{1}{C_{c}}\Big[\, u_2\dot m_0 c_p\,(T_{\mathrm{in,c}}-x_3) + UA_{\mathrm{ex}}\,(x_2-x_3)\Big].
\end{align}
These equations are derived from the energy balances. The first row of the state-space equation printed in the source paper (its Eq.~16) contains a sign inconsistency and is not copied; the second and third rows coincide with the paper.

\subsection{Cell voltage}
An Ulleberg-type correlation \cite{ulleberg2003} is used, with $\theta=T-273.15$ in $^\circ$C and active area $A$ in m$^2$:
\begin{equation}
\Vcell = V_{\mathrm{rev}} + \frac{r_1+r_2\theta}{A}\,I + \big(s_0+s_1\theta+s_2\theta^2\big)\ln\!\Big(\frac{t_1+t_2/\theta+t_3/\theta^2}{A}\,I+1\Big).
\end{equation}
The paper prints the logarithm argument with $t_2T+t_3T^2$; the original form $t_2/\theta+t_3/\theta^2$ is used here.

\subsection{Parameters and calibration}
Table~\ref{tab:params} lists the parameters. The values marked ``Paper'' are from the source paper. The paper does not give the active area, loss area, flow scale, chiller temperature or voltage constants, so these are assumed. $A_s$ and $T_{\mathrm{in,c}}$ were then calibrated (root finding) so that at $I=40$~A, $u_1=0.7$, $u_2=0.5$ the steady state matches the operating point of the source paper: $\Tout=320.0$~K, $\Tin=315.0$~K, $\Tc=312.9$~K, $\dT=5.0$~K and $\Vcell=1.870$~V. Nothing was tuned using AI results.

\begin{table}[H]
\centering\small
\caption{Parameters of the first-principles model.}\label{tab:params}
\begin{tabular}{@{}p{2.6cm}p{3.9cm}p{6.2cm}p{1.6cm}@{}}
\toprule
Symbol & Meaning & Value & Source \\
\midrule
$C_{\mathrm{th}}$, $C_h$, $C_c$ & Thermal capacities & $1.20\times10^{5}$, $1.46\times10^{5}$, $2.30\times10^{4}$ J\,K$^{-1}$ & Paper \\
$UA_{\mathrm{ex}}$ & Heat-exchanger conductance & 140 W\,K$^{-1}$ & Paper \\
$T_{\mathrm{amb}}$, $h$ & Ambient temperature, loss coefficient & 300 K, 3.8 W\,m$^{-2}$\,K$^{-1}$ & Paper \\
$n_{\mathrm{cell}}$, $V_{\mathrm{th}}$ & Cells, thermoneutral voltage & 26, 1.48 V & Paper \\
$c_p$, $V_{\mathrm{rev}}$ & Water specific heat, reversible voltage & 4180 J\,kg$^{-1}$\,K$^{-1}$, 1.23 V & Standard \\
$A$ & Active cell area & 0.25 m$^{2}$ & Assumed \\
$\dot m_0$ & Mass flow at $u=1$ & 0.02 kg\,s$^{-1}$ & Assumed \\
$A_s$ & Effective loss area & 1.48 m$^{2}$ & Calibrated \\
$T_{\mathrm{in,c}}$ & Chiller supply temperature & 305.9 K & Calibrated \\
$r_1,r_2$ & Ohmic terms & $4.45153\times10^{-5}$ $\Omega$\,m$^2$, $6.88874\times10^{-9}$ $\Omega$\,m$^2$\,$^\circ$C$^{-1}$ & Assumed \\
$s_0$ ($s_1=s_2=0$) & Overvoltage coefficient & 0.33824 V & Assumed \\
$t_1,t_2,t_3$ & Overvoltage constants & $-0.01539$, 2.00181, 15.24178 (m$^2$\,A$^{-1}$ with $^\circ$C powers) & Assumed \\
\bottomrule
\end{tabular}
\end{table}

\textbf{Important caveat.} The voltage constants are the commonly cited alkaline-electrolyzer set of Ulleberg \cite{ulleberg2003}. They were recalled from the literature and could not be verified against the paper text, and they describe an alkaline, not a PEM, cell. They serve only as a smooth, physically shaped $\Vcell(T,I)$ correlation. With them, $\Vcell$ lies between about 1.5 and 1.9~V over the operating range. All absolute temperature levels in this report therefore have no experimental backing.

\subsection{Numerical solution and verification}
The ODEs are integrated with fixed-step RK4 ($\Delta t=1$~s). Over the test trajectories, the maximum absolute difference to an adaptive integrator (SciPy \texttt{solve\_ivp}, DOP853, tolerances $10^{-11}$) is $2.9\times10^{-9}$~K. For fixed $(I,u_1,u_2)$ the steady state solves $f(\bm x^\ast,\bm u,I)=0$. The root finder starts from a 3000~s transient. A solution is accepted only if the residual norm is below $10^{-8}$, all Jacobian eigenvalues have negative real part, and all temperatures lie in $[295,360]$~K.

% ==================================================================
\section{Inputs and Outputs}
\begin{table}[H]
\centering\small
\caption{Variables, roles and ranges used for data generation (SS = steady state, DYN = dynamic).}\label{tab:vars}
\begin{tabular}{lllll}
\toprule
Variable & Description & Role & Range & Used in \\
\midrule
$I$ & Stack current [A] & Disturbance & $[5,50]$ & SS, DYN \\
$u_1$ & Hot-side valve opening & Manipulated input & $[0.3,1]$ & SS, DYN \\
$u_2$ & Cold-side valve opening & Manipulated input & $[0.05,1]$ & SS, DYN \\
$\Tout,\Tin,\Tc$ & Temperatures [K] & States / outputs & about $[304,340]$ & SS, DYN \\
$\Vcell$ & Cell voltage [V] & Output (efficiency) & $[1.51,1.92]$ & SS \\
$\dT=\Tout-\Tin$ & Temperature difference & Derived output & -- & SS, DYN \\
\bottomrule
\end{tabular}
\end{table}
The \textbf{steady-state model} is $\bm y=g(I,u_1,u_2)$ with $\bm y=[\Tout,\Tin,\Tc,\Vcell]^{\top}$. The \textbf{dynamic model} is a one-step predictor $\bm x_{k+1}=F(\bm x_k,\bm u_k,I_k)$ with $\Delta t=1$~s.

\textbf{Justification.} $u_1$ and $u_2$ are the only actuators and $I$ is the dominant disturbance, so these three variables determine the response once the constants are fixed. $\Tin$ and $\dT$ are the quantities to control in Stage~3; all three temperatures are needed for a Markovian dynamic model. $\Vcell$ is an extra steady-state output because it sets heat generation and efficiency, which Stage~2 needs. $T_{\mathrm{amb}}$ and $T_{\mathrm{in,c}}$ are fixed to keep the input space low-dimensional. The lower valve limits ($u_1\ge0.3$, $u_2\ge0.05$) keep both loops circulating; with $u_2\to0$ the cold loop removes no heat. The correlation matrix in Fig.~\ref{fig:sscorr} supports the selection: the current is strongly correlated with all outputs ($r=0.72$ to $0.93$), and $u_2$ is the main cooling lever ($r=-0.52$ with $\Tc$).

% ==================================================================
\section{Dataset Generation}
\subsection{Steady state}
$N=2000$ input triples are drawn by Latin hypercube sampling \cite{mckay1979} over the ranges of Table~\ref{tab:vars} and the steady state is solved for each. All 2000 solutions passed the acceptance checks, so none were dropped (Table~\ref{tab:data}). Figures~\ref{fig:sssampling} and~\ref{fig:ssdist} show the input coverage and the output distributions.

\fig{fig_ss_sampling}{0.85}{Latin-hypercube design in $(I,u_1,u_2)$: pairwise scatter of retained samples (no samples were dropped).}{fig:sssampling}
\fig{fig_ss_outputs_hist}{0.85}{Distributions of the steady-state outputs and of $\dT$.}{fig:ssdist}

\subsection{Dynamic}
Thirty trajectories of 2000~s ($\Delta t=1$~s, 60\,000 samples) are generated by RK4 under random multi-level excitation.
\begin{itemize}
  \item $I$, $u_1$, $u_2$ are piecewise constant with independent random hold times in $[20,100]$~s and uniform random levels. In half of the trajectories, step amplitudes are limited to $\pm20\%$ of the range around a random baseline.
  \item The initial state is the steady state of a random input triple plus Gaussian noise ($\sigma=1$~K).
  \item The split is by trajectory: 20 training, 5 validation, 5 test (Table~\ref{tab:data}).
  \item A noisy copy of the test set adds $N(0,0.05^2)$~K noise to the states; it is used only to test robustness.
\end{itemize}
Figures~\ref{fig:dyninputs} and~\ref{fig:dynresp} show an example excitation and response.

\begin{table}[H]
\centering\small
\caption{Dataset sizes and splits (dynamic: number of trajectories, one-step samples in brackets).}\label{tab:data}
\tab{dataset_summary}
\end{table}

\fig{fig_dyn_inputs}{0.85}{Example random multi-level excitation of $I$, $u_1$ and $u_2$.}{fig:dyninputs}
\fig{fig_dyn_response}{0.85}{Simulated temperature response to the excitation of Fig.~\ref{fig:dyninputs}.}{fig:dynresp}

\FloatBarrier
% ==================================================================
\section{Data Preprocessing}
\begin{enumerate}
  \item \textbf{Physical filtering.} Unstable or out-of-range equilibria are removed (none occurred). Because the data are simulated, no statistical outlier removal is applied.
  \item \textbf{Exploratory analysis.} Pairwise scatter plots and a correlation matrix (Fig.~\ref{fig:sscorr}).
  \item \textbf{Scaling.} Inputs and outputs are scaled to $[-1,1]$ (temperatures $(x-325)/25$, current $(I-27.5)/22.5$, valves by their mid-range and half-range, $\Vcell$ by its training minimum and maximum). Scaling constants come from the training data only.
  \item \textbf{Increment targets.} The temperature changes only about $10^{-3}$~K per second, so predicting $\bm x_{k+1}$ directly would reduce to learning the identity. The dynamic models therefore predict $\Delta\bm x_k=\bm x_{k+1}-\bm x_k$, scaled by its training standard deviation per state (0.0020, 0.0013 and 0.0067~K).
  \item \textbf{Lags and windows.} ARX and NARX use $\bm\xi_k=[\bm x_k,\bm x_{k-1},\bm u_k,\bm u_{k-1},I_k,I_{k-1}]$ (scaled); the LSTM uses windows of 10 steps with stride one. The first samples are edge-padded.
  \item \textbf{Splitting without leakage.} Steady-state data are split randomly 70/15/15~\%. Dynamic data are split by trajectory, never by rows. The test set is used only after all hyperparameters are chosen on the validation set.
\end{enumerate}

\fig{fig_ss_corr}{0.6}{Pearson correlation between steady-state inputs (rows) and outputs (columns).}{fig:sscorr}

\FloatBarrier
% ==================================================================
\section{AI Models}
All neural models are written in PyTorch and trained with Adam \cite{adam2015} on the MSE of the scaled targets, with early stopping on the validation loss. Hyperparameters are selected from small grids on the validation set (Appendix~B); the stochastic models are trained with three seeds and reported as mean $\pm$ standard deviation.

\subsection{Steady-state models}
\textbf{MLR} fits $\hat y_j=\beta_{0,j}+\sum_m\beta_{m,j}z_m$ with $\bm z=(I,u_1,u_2)$ by least squares. \textbf{Polynomial ridge} uses quadratic features and ridge regularization with $\lambda$ chosen on the validation set. \textbf{MLP}: tanh hidden layers and a linear output; grid over 2 or 3 layers, width 32 or 64 and learning rate $10^{-3}$ or $3\times10^{-3}$. \textbf{TSK / ANFIS-type} \cite{takagi1985,jang1993}: two Gaussian membership functions per input give $R=8$ rules,
\begin{equation}
\text{IF } I \text{ is } A_{r1}\text{ AND }u_1\text{ is }A_{r2}\text{ AND }u_2\text{ is }A_{r3}\ \text{THEN}\ y_r=\bm p_r^{\top}\bm z+q_r,
\end{equation}
with firing strengths $w_r=\prod_m\mu_{rm}(z_m)$ and output $\hat y=\sum_r\bar w_r y_r$, $\bar w_r=w_r/\sum_s w_s$. The consequents $(\bm p_r,q_r)$ are fitted by least squares and the premise parameters (centres, widths) by Adam, alternating five times.

\subsection{Dynamic models}
\textbf{ARX} regresses the increment on the lagged regressors, $\Delta\bm x_k=\bm\Theta\bm\xi_k+\bm b$, by ridge least squares. \textbf{NARX-MLP} feeds the same regressors to a tanh MLP with two hidden layers. \textbf{LSTM} \cite{hochreiter1997} is one recurrent layer with a linear head; it reads a window of the last 10 steps of $(\bm x,\bm u,I)$ and predicts $\Delta\bm x_k$. It is trained with teacher forcing and gradient clipping at 1.0, and in free run it is fed its own predictions.

\textbf{PINC (self-loop PINN).} The network learns the one-step transition map \cite{antonelo2024,zerrougui2026}. It takes the initial state, the held inputs and a normalized time $\tau\in[0,1]$ inside the step. With the hard initial condition used here,
\begin{equation}
\hat{\bm x}(\tau)=\bm x_k+\tau\,\bm\sigma_{\Delta}\odot\mathcal N_\theta(\bm x_k,\bm u_k,I_k,\tau),\qquad \hat{\bm x}_{k+1}=\hat{\bm x}(1),
\end{equation}
so $\hat{\bm x}(0)=\bm x_k$ holds exactly ($\bm\sigma_\Delta$ is the fixed standard deviation of the increments). Predictions are chained by feeding $\hat{\bm x}_{k+1}$ back (Fig.~\ref{fig:pincarch}). The loss is the ODE residual at $N_c=8$ random collocation points per sample,
\begin{equation}
\mathcal L_{\mathrm{phys}}=\mathrm{mean}\Big\|\Big(\tfrac{1}{\Delta t}\partial_\tau\hat{\bm x}-f(\hat{\bm x},\bm u_k,I_k)\Big)\oslash\bm\sigma_{\dot x}\Big\|_2^2,
\end{equation}
with $\partial_\tau$ from automatic differentiation and $\bm\sigma_{\dot x}$ the standard deviation of $f$ over the training states. \textbf{No next-state targets from the simulator are used.} Training points are the states and inputs of the training trajectories plus 30\,\% uniform random points in the operating box. The network has four hidden layers of 64 tanh neurons; optimization is Adam (5000 iterations, learning rate $10^{-3}\to10^{-4}$) followed by L-BFGS \cite{liu1989} (500 iterations). A soft-IC variant (paper formulation, adaptive loss weights) is available in the code but was not used for the results.

\fig{fig_pinc_architecture}{0.85}{PINC: the network maps $(\bm x_k,\bm u_k,I_k,\tau)$ to the state inside the step; the end state is fed back as the next initial state; the physics loss uses automatic differentiation.}{fig:pincarch}

\FloatBarrier
% ==================================================================
\section{Performance Metrics}
For $N$ test points with true values $y_n$ and predictions $\hat y_n$:
\begin{align}
\mathrm{RMSE}&=\sqrt{\tfrac1N\textstyle\sum_n(y_n-\hat y_n)^2}, &
\mathrm{MAE}&=\tfrac1N\textstyle\sum_n|y_n-\hat y_n|,\\
R^2&=1-\frac{\sum_n(y_n-\hat y_n)^2}{\sum_n(y_n-\bar y)^2}, &
\mathrm{NRMSE}&=\frac{\mathrm{RMSE}}{y_{\max}-y_{\min}},
\end{align}
and the maximum absolute error. RMSE is the main measure because the later stages need absolute temperature accuracy. Errors are reported per output and never summed across outputs of different units. Dynamic models are evaluated in three ways: (i) one-step error with teacher forcing, plus $R^2_\Delta$ computed on the increments $\Delta\bm x$; (ii) free-run rollout, starting from the true initial state of each test trajectory, driven by the true inputs and fed back with its own predictions over 2000~s. The rollout RMSE at horizon $H$ is pooled over the three temperatures and over steps $1\ldots H$; (iii) cost (parameters, training time, inference time per 1000 steps relative to RK4). Robustness is checked on the noisy test set.

% ==================================================================
\section{Results and Discussion}
\subsection{Steady-state models}
Table~\ref{tab:ssmetrics} gives test-set results. The MLP is clearly the best model: RMSE $0.042$, $0.036$ and $0.043$~K for the three temperatures and $4.4\times10^{-4}$~V for $\Vcell$ (NRMSE about 0.13\,\% of the range), mean $R^2\approx1.000$ and a maximum temperature error of 0.26~K (Fig.~\ref{fig:ssparity}). The linear model has a mean $R^2$ of 0.85 and errors up to 11.7~K; its residuals show a strong curved pattern (Fig.~\ref{fig:ssresid}), which means that the response is nonlinear in the inputs. Quadratic features remove most of this (RMSE about 1~K) and the TSK model reduces it further (about 0.4~K), but neither removes the pattern entirely; the MLP residuals are small and unstructured (within about $\pm0.5$\,\% of the range). The MLP learning curve (Fig.~\ref{fig:sslearn}) shows no overfitting: validation loss follows training loss. The selected model was still improving slowly when training stopped at the epoch limit of 1000 (best epoch 971), so a longer budget could improve it slightly.

\begin{table}[H]
\centering\small
\caption{Steady-state test performance (RMSE in K for temperatures, V for $\Vcell$; $R^2$ averaged over outputs; max $|e_T|$ over the three temperatures). The MLP row is mean $\pm$ std over 3 seeds.}\label{tab:ssmetrics}
\resizebox{\textwidth}{!}{\tab{ss_metrics}}
\end{table}

\fig{fig_ss_parity}{0.95}{Parity plots on the test set (rows: outputs; columns: models).}{fig:ssparity}
\fig{fig_ss_residuals}{0.95}{Residuals (percent of output range) versus predicted value and residual histograms on the test set.}{fig:ssresid}
\fig{fig_ss_learning_curve}{0.6}{Training and validation loss of the selected steady-state MLP.}{fig:sslearn}

\subsection{Dynamic models: one-step prediction}
All models predict one step ahead with errors of the order of $10^{-5}$ to $10^{-4}$~K (Table~\ref{tab:dynone}, Fig.~\ref{fig:dynone}). Since the states change by only a few millikelvin per step, the increment $R^2_\Delta$ is more informative: NARX-MLP 0.9994, PINC 0.9984, ARX 0.9956 and LSTM 0.9895. The PINC reaches almost the accuracy of the supervised NARX-MLP without ever seeing a simulated next state.

\begin{table}[H]
\centering\small
\caption{One-step (teacher-forced) accuracy on the test trajectories, mean $\pm$ std over seeds (ARX has no seed variance).}\label{tab:dynone}
\resizebox{\textwidth}{!}{\tab{dyn_onestep_metrics}}
\end{table}

\fig{fig_dyn_onestep}{0.95}{One-step predictions versus the simulator on one test trajectory (left) and one-step errors (right).}{fig:dynone}

\subsection{Dynamic models: free-run rollout}
Feeding predictions back exposes drift (Table~\ref{tab:dynroll}, Figs.~\ref{fig:dynroll} and~\ref{fig:rollerr}). Over the full 2000~s the pooled RMSE is 0.012~K for NARX-MLP, 0.023~K for PINC, 0.038~K for LSTM and 0.46~K for ARX.
\begin{itemize}
  \item \textbf{NARX-MLP} is best over the full horizon; its error rises to about 0.012~K within 100~s and then stays flat (bounded).
  \item \textbf{PINC} is the most accurate for short horizons (0.0036~K at $H=100$ versus 0.0117~K for NARX-MLP) but drifts slowly, roughly like $t^{0.5}$ between 100 and 2000~s, to 0.023~K. The physics loss does not guarantee that the one-step error stays unbiased when accumulated.
  \item \textbf{LSTM} reaches 0.056~K by $H=100$ and then settles near 0.04~K (bounded).
  \item \textbf{ARX} drifts steadily (0.10~K at $H=100$, 0.46~K at $H=2000$); a linear model cannot represent the nonlinear heat generation and cooling terms.
\end{itemize}
All rollouts remain stable (no divergence). The open-loop step responses in Fig.~\ref{fig:stepresp}, which matter for Stage~3, are followed closely by NARX-MLP, PINC and LSTM (500~s RMSE below about 0.05~K for all three states), whereas ARX deviates by up to 0.35~K for a step in $u_2$.

\begin{table}[H]
\centering\small
\caption{Free-run rollout RMSE on the test trajectories (K, mean $\pm$ std over seeds). $H$ columns are pooled over the three temperatures; the last three columns are per state over the full 2000~s.}\label{tab:dynroll}
\resizebox{\textwidth}{!}{\tab{dyn_rollout_metrics}}
\end{table}

\fig{fig_dyn_rollout}{0.95}{Free-run rollout of all dynamic models versus the simulator on a test trajectory.}{fig:dynroll}
\fig{fig_dyn_rollout_error}{0.75}{Rollout RMSE versus horizon (mean over seeds, shaded band: standard deviation).}{fig:rollerr}
\fig{fig_dyn_step_response}{0.95}{Open-loop step responses from a common nominal state ($I=30$~A, $u_1=u_2=0.6$): step in $I$ (30 to 45~A), in $u_1$ and in $u_2$ (0.6 to 0.9).}{fig:stepresp}

\subsection{PINC training}
Figure~\ref{fig:pincloss} shows the physics loss (normalised residual squared). In the Adam stage it falls from about 1 to about $10^{-3}$; L-BFGS reduces it further to $2\times10^{-5}$ (about 4.5 orders of magnitude overall). Training took about 13 min per seed (776~s on average) on the CPU. On the test states, the mean squared normalised ODE residual of the learned map is about $1.0\times10^{-5}$, i.e.\ an RMS residual of about $1.7\times10^{-4}$~K\,s$^{-1}$, an indicator of physical consistency that does not use simulated trajectories.

\fig{fig_pinc_loss}{0.9}{PINC loss history: Adam stage (left) and L-BFGS stage (right, per function evaluation).}{fig:pincloss}

\subsection{Computational cost}
\begin{table}[H]
\centering\small
\caption{Parameters, training time and inference time per 1000 rollout steps (median of repeats, one trajectory) compared with RK4.}\label{tab:cost}
\tab{dyn_cost}
\end{table}
For a single three-state, non-stiff simulation the surrogates give no speed-up: NARX-MLP takes about the same time as RK4 (1.04 times), PINC 1.4 times, LSTM 4.4 times, and only the linear ARX is faster (0.5 times). The benefit of the surrogates is that they are differentiable and can be batched, which matters for Stages~2 and~3. The PINC costs far more to train (776~s versus 24~s for NARX-MLP), since each step involves second-order derivatives through the network.

\subsection{Robustness to measurement noise}
Models trained on clean data are tested with noisy test states ($\sigma_n=0.05$~K), where the rollout starts from the noisy measured state (Table~\ref{tab:noise}). The effect on free-run rollouts is moderate: ARX is unchanged (its error is dominated by bias), NARX-MLP degrades 2.5 times (0.012 to 0.030~K), PINC 1.4 times and LSTM 1.3 times. Noise matters much more for one-step prediction from noisy measured states, where all models reach 0.05 to 0.11~K, because the noise (0.05~K) is far larger than the per-step increments (about 0.002~K). This is why the models would need filtering or state estimation in a closed loop with real sensors.

\begin{table}[H]
\centering\small
\caption{Rollout RMSE (K) on clean and noisy test data.}\label{tab:noise}
\tab{dyn_noise_metrics}
\end{table}

\subsection{Discussion and selected models}
\begin{itemize}
  \item \textbf{Steady state.} The input--output map is nonlinear enough that a linear model is inadequate (mean $R^2=0.85$). The MLP is accurate to about 0.04~K and is selected as the Stage~2 surrogate: it is smooth, differentiable and also predicts $\Vcell$.
  \item \textbf{Dynamics.} NARX-MLP is the most accurate over long horizons and is selected as the Stage~3 prediction model. The PINC is a strong alternative, since it is the best at short horizons, needs no simulated next states, and its physical consistency can be checked, but it trains much more slowly and accumulates a small drift.
  \item \textbf{Role of physics.} The PINC matches the supervised models closely using only the governing equations. It is a data-efficient option when simulated or measured next states are scarce.
  \item \textbf{Limitations.} All results are surrogate-versus-simulator comparisons on one lumped model; there are no experimental data. Several parameters are assumed, and the voltage constants come from an alkaline set that was not verified. Rollouts are tested up to 2000~s, and the models are valid only inside the sampled operating box. A single dataset realization was used; only the neural-network seeds vary (three seeds).
\end{itemize}

% ==================================================================
\section{Conclusions}
A three-state thermal model of a PEM electrolyzer was implemented, verified against an adaptive solver (difference $2.9\times10^{-9}$~K) and used to generate steady-state and dynamic datasets. Steady-state surrogates (MLR, polynomial ridge, MLP, TSK/ANFIS-type) and dynamic surrogates (ARX, NARX-MLP, LSTM, PINC) were trained with a leakage-free protocol and compared on held-out data. The MLP is the best steady-state model (RMSE 0.04~K, $R^2\approx1$), and the NARX-MLP the best dynamic model (2000~s rollout RMSE 0.012~K); the PINC reaches comparable accuracy without next-state data. The MLP and NARX-MLP are carried forward to Stage~2 (optimization) and Stage~3 (control).

% ==================================================================
\begin{thebibliography}{99}
\bibitem{zerrougui2026} I.~Zerrougui, Z.~Li, D.~Hissel, ``Self-loop physics-informed neural network for model predictive control of PEM electrolyzers,'' \emph{Journal of Process Control}, vol.~161, 103707, 2026.
\bibitem{antonelo2024} E.~A.~Antonelo, E.~Camponogara, L.~O.~Seman, J.~P.~Jordanou, E.~R.~de Souza, J.~F.~H\"ubner, ``Physics-informed neural nets for control of dynamical systems,'' \emph{Neurocomputing}, vol.~579, 127419, 2024.
\bibitem{ulleberg2003} {\O}.~Ulleberg, ``Modeling of advanced alkaline electrolyzers: a system simulation approach,'' \emph{Int.\ J.\ Hydrogen Energy}, vol.~28, no.~1, pp.~21--33, 2003.
\bibitem{carmo2013} M.~Carmo, D.~L.~Fritz, J.~Mergel, D.~Stolten, ``A comprehensive review on PEM water electrolysis,'' \emph{Int.\ J.\ Hydrogen Energy}, vol.~38, no.~12, pp.~4901--4934, 2013.
\bibitem{mckay1979} M.~D.~McKay, R.~J.~Beckman, W.~J.~Conover, ``A comparison of three methods for selecting values of input variables in the analysis of output from a computer code,'' \emph{Technometrics}, vol.~21, no.~2, pp.~239--245, 1979.
\bibitem{adam2015} D.~P.~Kingma, J.~Ba, ``Adam: A method for stochastic optimization,'' in \emph{Proc.\ ICLR}, 2015.
\bibitem{liu1989} D.~C.~Liu, J.~Nocedal, ``On the limited memory BFGS method for large scale optimization,'' \emph{Math.\ Programming}, vol.~45, pp.~503--528, 1989.
\bibitem{hochreiter1997} S.~Hochreiter, J.~Schmidhuber, ``Long short-term memory,'' \emph{Neural Computation}, vol.~9, no.~8, pp.~1735--1780, 1997.
\bibitem{jang1993} J.-S.~R.~Jang, ``ANFIS: Adaptive-network-based fuzzy inference system,'' \emph{IEEE Trans.\ Syst., Man, Cybern.}, vol.~23, no.~3, pp.~665--685, 1993.
\bibitem{takagi1985} T.~Takagi, M.~Sugeno, ``Fuzzy identification of systems and its applications to modeling and control,'' \emph{IEEE Trans.\ Syst., Man, Cybern.}, vol.~15, no.~1, pp.~116--132, 1985.
\bibitem{lectures} V.~S.~Kumar, ``Applications of AI in Chemical Engineering --- AI in Process Modeling,'' lecture slides, CHL~4380, IIT Jodhpur.
\end{thebibliography}

\clearpage
% ==================================================================
\appendix
\section{Reproducibility}
\begin{itemize}
  \item \textbf{Run.} \texttt{pip install -r requirements.txt}, then \texttt{python scripts/run\_all.py} (about 45--50 min; \texttt{--fast} is a smoke test). \texttt{python scripts/05\_make\_diagrams.py} redraws the three block diagrams.
  \item \textbf{Settings.} All sizes, ranges, seeds and grids are in \texttt{config.yaml}; process parameters with their sources are in \texttt{params.yaml}. Seeds are fixed; three seeds are used for the neural models.
  \item \textbf{Report build.} This file sits in the repository root. Compile with \texttt{pdflatex} twice; figures come from \texttt{figures/} and table bodies from \texttt{results/tables/}, both written by the code.
  \item \textbf{Software.} Python 3.12.4, NumPy 2.5.3, SciPy 1.18.1, scikit-learn 1.9.1, PyTorch 2.14.1 (CPU, 2 threads). Hardware: Intel Core 7 150U laptop CPU.
\end{itemize}

\section{Selected Hyperparameters}
\begin{table}[H]
\centering\footnotesize
\caption{Hyperparameters selected on the validation set (dynamic networks by 100-step validation rollout RMSE).}\label{tab:hyper}
\tab{hyperparams}
\end{table}

\end{document}
